\section{结论}
\label{sec:conclusion}

本文针对大型转动设备轴系对中中“空间任意点到激光基准轴垂直距离”这一在现有商业
产品与文献中尚未被直接解决的测量问题，提出了 DCPAM 测量模型。该模型以一条激光束
替代传统钢丝找中法中的张紧钢丝，消除了钢丝下坠挠度带来的系统偏差，并免除了每次
测量时反复架设钢丝所引发的工期干扰与安全隐患。方法上，通过前、后两个相机模块
拍摄激光在毛玻璃接收屏上的光斑，经反投影、基于五圆点 PnP 的位姿对齐、设备坐标系
下的镜面虚像重建与叉积点线距离计算，把激光基准轴还原为设备坐标系下可解析的三维
直线，从而对任意目标点直接求解其到该轴的垂距。

围绕重复精度与误差传播的实验与分析表明，系统的主导误差来源位于测量端的圆心提取
像素噪声，并经 $Z$ 方向的光路杠杆放大，而设备几何参数已接近物理真值、并非方差的
主要杠杆；改进圆心提取显著降低了重复测量的组内标准差。当前重复精度与
$\pm15$--$50\,\mu\mathrm{m}$ 工程目标之间的差距及其定量结论仍在补充实验中完善。

后续工作将从三方面推进。其一，在测量端进一步压低像素噪声，包括用卷积神经网络直接
从原始图像回归光斑圆心以替代部分手工特征提取，以及改用更高信噪比的相机与激光源。
其二，建立测杆倾斜偏差 $\beta$ 与测量精度的关系模型，给出在指定精度下的倾斜度
舍弃阈值，并通过光路与靶点几何设计降低 $Z$ 向杠杆敏感度。其三，完善全量关节臂
扫描标定与标定一致性自检，进一步巩固几何真值并支撑现场大批量采样点的对中作业。
